\documentclass[11pt,letterpaper]{article}
\usepackage{../../Common/course_notes}
\usepackage{../../Common/course_macros}

\begin{document}
\setlength{\abovedisplayskip}{6pt}
\setlength{\belowdisplayskip}{6pt}
\setlength{\textfloatsep}{10pt}
\setlength{\intextsep}{8pt}
\color{blackmain}

\begin{center}
{\LARGE\textbf{Geometría y Física I: Mecánica Clásica y Formas de Cartan}}\\[0.25cm]
{\Large\textcolor{orangeacc}{\textbf{Semana 2: De la Crisis Cartesiana al Arsenal de Índices}}}\\[0.15cm]
{\large\textbf{Construyendo la notación de Einstein, $\delta_{ij}$ y $\epsilon_{ijk}$ desde la geometría}}\\[0.3cm]
{\large\textcolor{graymuted}{Notas de Estudio --- Módulo 0: El Arsenal Indicial y Geométrico}}\\[0.2cm]
\end{center}

\vspace{0.2cm}

% =====================================================================
\manualsection[1. La Crisis Cartesiana]{1. La Crisis del Método Tradicional: ¿Por qué necesitamos otra herramienta?}

En los cursos elementales de cálculo multivariable aprendimos a manipular vectores en $\R^3$ separando sus componentes a lo largo de tres ejes fijos $(x, y, z)$ con vectores unitarios $\{\hat{i}, \hat{j}, \hat{k}\}$. Para calcular un producto cruz, solemos recurrir al determinante mnemotécnico:
\[
\vec{A} \times \vec{B} = \det \begin{pmatrix} \hat{i} & \hat{j} & \hat{k} \\ A_x & A_y & A_z \\ B_x & B_y & B_z \end{pmatrix}
= (A_y B_z - A_z B_y)\hat{i} + (A_z B_x - A_x B_z)\hat{j} + (A_x B_y - A_y B_x)\hat{k}.
\]

¿Qué ocurre cuando un problema físico real exige encadenar dos productos cruz? En electromagnetismo y dinámica de cuerpos rígidos aparece constantemente el \textbf{doble producto vectorial} $\vec{A} \times (\vec{B} \times \vec{C})$. Intentemos calcular únicamente su primera componente (la componente $x$) usando el método tradicional:

\begin{examplebox}{Experimento Preliminar: Calculando $[\vec{A}\times(\vec{B}\times\vec{C})]_x$ a mano}
Llamemos $\vec{D} = \vec{B} \times \vec{C}$. Sabemos que sus componentes $y$ y $z$ son:
\[
D_y = B_z C_x - B_x C_z, \qquad D_z = B_x C_y - B_y C_x.
\]
Ahora calculamos la componente $x$ de $\vec{A} \times \vec{D}$:
\begin{align*}
[\vec{A} \times (\vec{B} \times \vec{C})]_x &= A_y D_z - A_z D_y \\
&= A_y(B_x C_y - B_y C_x) - A_z(B_z C_x - B_x C_z) \\
&= B_x(A_y C_y + A_z C_z) - C_x(A_y B_y + A_z B_z).
\end{align*}
Observa la expresión entre paréntesis: $A_y C_y + A_z C_z$ \textit{casi} es el producto punto $\vec{A}\cdot\vec{C}$, pero le falta el término $A_x C_x$. Para completar el producto punto, nos vemos obligados a \textbf{sumar y restar artificialmente} el término $A_x B_x C_x$:
\begin{align*}
[\vec{A} \times (\vec{B} \times \vec{C})]_x &= B_x(A_x C_x + A_y C_y + A_z C_z) - C_x(A_x B_x + A_y B_y + A_z B_z) \\
&= B_x(\vec{A}\cdot\vec{C}) - C_x(\vec{A}\cdot\vec{B}).
\end{align*}
Finalmente, habría que repetir exactamente este mismo malabarismo algebraico dos veces más para las componentes $y$ y $z$.
\end{examplebox}

\begin{intuitionbox}{Diagnóstico de la Crisis}
¿Por qué este cálculo fue tan torpe y requirió un «truco» artificial de sumar y restar $A_x B_x C_x$?
\begin{enumerate}
    \item \textbf{Repetición ciega:} Las componentes $x, y, z$ obedecen exactamente la misma ley geométrica, pero la notación tradicional nos obliga a tratarlas por separado.
    \item \textbf{Pérdida de estructura:} Al expandir prematuramente el determinante, destruimos la propiedad geométrica fundamental del producto cruz (la antisimetría y el volumen) y nos quedamos con una sopa de letras.
\end{enumerate}
\textbf{Meta de la sesión:} Construir desde cero un lenguaje donde las tres componentes se manipulen simultáneamente y donde la geometría de la base ortonormal haga el trabajo pesado por nosotros.
\end{intuitionbox}

% =====================================================================
\manualsection[2. Construyendo $\delta_{ij}$]{2. Extracción de Componentes: Nacimiento de Einstein y $\delta_{ij}$}

Para tratar los tres ejes en pie de igualdad, rebautizamos las coordenadas $(x, y, z) \to (x_1, x_2, x_3)$ y los vectores unitarios de la base canónica como:
\[
\hat{e}_1 \equiv \hat{i}, \qquad \hat{e}_2 \equiv \hat{j}, \qquad \hat{e}_3 \equiv \hat{k}.
\]
Cualquier vector $\vec{A} \in \R^3$ se escribe como la combinación lineal:
\[
\vec{A} = A_1 \hat{e}_1 + A_2 \hat{e}_2 + A_3 \hat{e}_3 = \sum_{j=1}^3 A_j \hat{e}_j.
\]

Hagámonos ahora una pregunta constructiva elemental: \textbf{Si conocemos el vector geométrico $\vec{A}$, ¿qué operación debemos hacerle para extraer únicamente su $i$-ésima componente $A_i$?}

Proyectamos $\vec{A}$ haciendo el producto escalar con el vector unitario $\hat{e}_i$:
\[
\hat{e}_i \cdot \vec{A} = \hat{e}_i \cdot \left(\sum_{j=1}^3 A_j \hat{e}_j\right) = \sum_{j=1}^3 A_j (\hat{e}_i \cdot \hat{e}_j).
\]
Observemos esta ecuación con detenimiento. Para que la suma de la derecha colapse y nos devuelva exclusivamente la componente $A_i$, toda la información geométrica reside en los nueve productos escalares $\hat{e}_i \cdot \hat{e}_j$ entre los vectores de la base. Como los ejes son unitarios y mutuamente perpendiculares:
\begin{itemize}
    \item Si $i = j$, el vector se proyecta sobre sí mismo: $\hat{e}_1\cdot\hat{e}_1 = \hat{e}_2\cdot\hat{e}_2 = \hat{e}_3\cdot\hat{e}_3 = 1$.
    \item Si $i \neq j$, los vectores están a $90^\circ$: $\hat{e}_1\cdot\hat{e}_2 = \hat{e}_2\cdot\hat{e}_3 = \hat{e}_3\cdot\hat{e}_1 = 0$.
\end{itemize}

Esta tabla de multiplicar geométrica merece un símbolo propio. Así nace la \term{delta de Kronecker}:

\begin{theorembox}{Construcción Geométrica de la Delta de Kronecker $\delta_{ij}$}
Definimos el símbolo $\delta_{ij}$ como el producto escalar entre los vectores unitarios de una base ortonormal:
\[
\mathbf{\delta_{ij} \equiv \hat{e}_i \cdot \hat{e}_j} = \begin{cases} 1 & \text{si } i = j, \\ 0 & \text{si } i \neq j. \end{cases}
\]
\end{theorembox}

Al mismo tiempo, notamos que el símbolo de sumatoria $\sum_{j=1}^3$ aparece siempre que un índice (aquí la letra $j$) se repite dos veces en el mismo término. Siguiendo a Einstein, suprimimos el símbolo $\sum$ estableciendo que \textbf{todo índice repetido dos veces en un monomio implica suma automática de $1$ a $3$}.

Con estas dos creaciones, la extracción de componentes se escribe de forma limpia:
\[
A_i = \hat{e}_i \cdot \vec{A} = \hat{e}_i \cdot (A_j \hat{e}_j) = A_j (\hat{e}_i \cdot \hat{e}_j) \implies \mathbf{\delta_{ij} A_j = A_i}.
\]
Y el producto escalar de dos vectores arbitrarios $\vec{A} = A_i \hat{e}_i$ y $\vec{B} = B_j \hat{e}_j$ se deduce sin memorizar fórmulas:
\[
\vec{A} \cdot \vec{B} = (A_i \hat{e}_i) \cdot (B_j \hat{e}_j) = A_i B_j (\hat{e}_i \cdot \hat{e}_j) = \delta_{ij} A_i B_j = \mathbf{A_i B_i}.
\]

\begin{infobox}[pureblack]{Higiene de Índices: Libres vs. Mudos}
\begin{itemize}
    \item \term{Índice libre:} Aparece \textbf{una sola vez} en cada término aditivo (como $i$ en $A_i = \delta_{ij} A_j$). Indica qué componente estamos mirando ($1, 2$ o $3$) y debe coincidir a ambos lados de la ecuación.
    \item \term{Índice mudo:} Aparece \textbf{exactamente dos veces} en un término (como $j$ en $\delta_{ij} A_j$). Está sumado internamente ($\delta_{i1}A_1 + \delta_{i2}A_2 + \delta_{i3}A_3$), por lo que podemos cambiarle la letra a voluntad: $\delta_{ij} A_j = \delta_{ik} A_k$.
    \item \textbf{Prohibición absoluta:} Un mismo índice \textbf{nunca} puede aparecer tres o más veces en un solo producto (como $A_i B_i C_i$), pues sería ambiguo qué par se suma.
\end{itemize}
\end{infobox}

% =====================================================================
\manualsection[3. Construyendo $\epsilon_{ijk}$]{3. El Producto Cruz sin Determinantes: Nacimiento de $\epsilon_{ijk}$}

Apliquemos exactamente la misma estrategia constructiva al producto cruz. En lugar de usar el determinante $3\times 3$, multipliquemos $\vec{A} = A_j \hat{e}_j$ y $\vec{B} = B_k \hat{e}_k$ usando únicamente la propiedad distributiva (bilinealidad):
\[
\vec{A} \times \vec{B} = (A_j \hat{e}_j) \times (B_k \hat{e}_k) = A_j B_k (\hat{e}_j \times \hat{e}_k).
\]
*(Nota que usamos dos letras distintas, $j$ y $k$, porque son dos sumas independientes).*

Ahora preguntamos: \textbf{¿Cómo extraemos la $i$-ésima componente del vector $\vec{A} \times \vec{B}$?} Igual que antes, proyectamos haciendo el producto punto con $\hat{e}_i$:
\[
(\vec{A} \times \vec{B})_i = \hat{e}_i \cdot (\vec{A} \times \vec{B}) = A_j B_k \underbrace{\big[\hat{e}_i \cdot (\hat{e}_j \times \hat{e}_k)\big]}_{\text{¿Qué es este número?}}.
\]

Detengámonos a admirar este resultado. Toda la estructura del producto cruz se ha concentrado en el número real $\hat{e}_i \cdot (\hat{e}_j \times \hat{e}_k)$. Del cálculo elemental sabemos que el triple producto escalar $\vec{u} \cdot (\vec{v} \times \vec{w})$ mide el **volumen con signo (orientado)** del paralelepípedo formado por los tres vectores. Bautizamos a este volumen elemental como el \term{símbolo de Levi-Civita} $\epsilon_{ijk}$:

\begin{theorembox}{Construcción Geométrica del Símbolo de Levi-Civita $\epsilon_{ijk}$}
Definimos $\epsilon_{ijk}$ como el volumen orientado del paralelepípedo generado por la terna de vectores base $(\hat{e}_i, \hat{e}_j, \hat{e}_k)$:
\[
\mathbf{\epsilon_{ijk} \equiv \hat{e}_i \cdot (\hat{e}_j \times \hat{e}_k)}.
\]
En consecuencia, la $i$-ésima componente del producto cruz es inmediatamente:
\[
\mathbf{(\vec{A} \times \vec{B})_i = \epsilon_{ijk} A_j B_k}.
\]
\end{theorembox}

A partir de esta definición geométrica como volumen, **no necesitamos memorizar las reglas de $\epsilon_{ijk}$; podemos deducirlas visualmente**:

\begin{enumerate}
    \item \textbf{¿Qué pasa si dos índices se repiten} (por ejemplo, $\epsilon_{112} = \hat{e}_1 \cdot (\hat{e}_1 \times \hat{e}_2)$)?\\
    Si dos aristas del paralelepípedo son el mismo vector, la caja queda completamente aplastada en un plano. Su volumen es cero:
    \[
    \mathbf{\epsilon_{iik} = \epsilon_{iji} = \epsilon_{jkk} = 0} \quad (\text{21 de las 27 componentes son cero}).
    \]
    \item \textbf{¿Qué pasa si los tres índices son $(1, 2, 3)$ en orden cíclico}?\\
    Los vectores $(\hat{e}_1, \hat{e}_2, \hat{e}_3)$ forman un cubo unitario con orientación de mano derecha ($\hat{e}_2 \times \hat{e}_3 = \hat{e}_1$). Su volumen es $+1$:
    \[
    \mathbf{\epsilon_{123} = \epsilon_{231} = \epsilon_{312} = +1}.
    \]
    \item \textbf{¿Qué pasa si intercambiamos dos vectores adyacentes}?\\
    Como $\hat{e}_k \times \hat{e}_j = -\hat{e}_j \times \hat{e}_k$, intercambiar cualquier par de índices invierte la orientación de la caja y cambia el signo del volumen:
    \[
    \mathbf{\epsilon_{132} = \epsilon_{213} = \epsilon_{321} = -1}.
    \]
\end{enumerate}

\begin{center}
\begin{tikzpicture}[scale=1.05]
    % Left: Unit Cube Volume
    \begin{scope}[xshift=-4cm, yshift=-0.8cm]
        \fill[orangebg] (0,0) -- (1.6,0) -- (2.2,0.6) -- (0.6,0.6) -- cycle;
        \fill[creambg] (0,0) -- (0,1.6) -- (0.6,2.2) -- (0.6,0.6) -- cycle;
        \draw[dashed, color=graymuted] (0.6,0.6) -- (2.2,0.6);
        \draw[dashed, color=graymuted] (0.6,0.6) -- (0.6,2.2);
        \draw[dashed, color=graymuted] (0,0) -- (0.6,0.6);
        \draw[line width=0.8pt, color=graymuted] (1.6,0) -- (2.2,0.6) -- (2.2,2.2) -- (1.6,1.6) -- cycle;
        \draw[line width=0.8pt, color=graymuted] (0,1.6) -- (0.6,2.2) -- (2.2,2.2) -- (1.6,1.6) -- cycle;
        \draw[line width=0.8pt, color=graymuted] (0,0) rectangle (1.6,1.6);
        
        \draw[-{Stealth[length=2.8mm]}, line width=1.6pt, color=navyaccent] (0,0) -- (1.6,0) node[below=2pt] {$\hat{e}_1$};
        \draw[-{Stealth[length=2.8mm]}, line width=1.6pt, color=navyaccent] (0,0) -- (0.6,0.6) node[above right=-1pt] {$\hat{e}_2$};
        \draw[-{Stealth[length=2.8mm]}, line width=1.6pt, color=orangeacc] (0,0) -- (0,1.6) node[left=2pt] {$\hat{e}_3$};
        \node[font=\small\bfseries, align=center] at (1.1, -0.7) {$\epsilon_{123} = \text{Vol}(\hat{e}_1,\hat{e}_2,\hat{e}_3) = +1$};
    \end{scope}

    % Right: Cyclic Clock
    \begin{scope}[xshift=2.5cm]
        \draw[line width=1.5pt, color=orangeacc] (0,0) circle (1.35cm);
        \node[circle, fill=blackmain, text=white, font=\bfseries, inner sep=3.5pt] at (90:1.35cm) {1};
        \node[circle, fill=blackmain, text=white, font=\bfseries, inner sep=3.5pt] at (330:1.35cm) {2};
        \node[circle, fill=blackmain, text=white, font=\bfseries, inner sep=3.5pt] at (210:1.35cm) {3};
        \draw[-{Stealth[length=2.5mm]}, line width=1.5pt, color=orangeacc] (65:1.65cm) arc (65:15:1.65cm) node[midway, right=2pt, font=\footnotesize\bfseries, color=orangeacc] {$+1$ (horario)};
        \draw[-{Stealth[length=2.5mm]}, line width=1.5pt, color=orangeacc] (305:1.65cm) arc (305:255:1.65cm);
        \draw[-{Stealth[length=2.5mm]}, line width=1.5pt, color=orangeacc] (185:1.65cm) arc (185:135:1.65cm);
        \node[align=center, font=\footnotesize\bfseries] at (0, 0) {Reloj de\\Orientación};
    \end{scope}
\end{tikzpicture}
\end{center}

% =====================================================================
\manualsection[4. Construyendo la Identidad $\epsilon-\delta$]{4. ¿De dónde sale la Identidad $\epsilon-\delta$? Construcción por Conteo}

Volvamos ahora a nuestro problema original: calcular $\vec{A} \times (\vec{B} \times \vec{C})$. En nuestra nueva notación, la $i$-ésima componente es:
\[
[\vec{A} \times (\vec{B} \times \vec{C})]_i = \epsilon_{ijk} A_j (\vec{B} \times \vec{C})_k = \epsilon_{ijk} A_j (\epsilon_{kmn} B_m C_n) = \mathbf{\epsilon_{kij}\epsilon_{kmn}} A_j B_m C_n,
\]
donde usamos la simetría cíclica $\epsilon_{ijk} = \epsilon_{kij}$ para colocar el índice mudo compartido $k$ al principio de ambos símbolos.

Vemos que la clave de todo el cálculo está en el producto contraído $\epsilon_{kij}\epsilon_{kmn}$ (o escribiendo $i$ como el índice mudo, $\mathbf{\epsilon_{ijk}\epsilon_{imn}}$). En lugar de aceptar una fórmula mágica, \textbf{construyámosla nosotros mismos analizando cuándo no es cero}:

\begin{intuitionbox}{Deducción Constructiva de $\epsilon_{ijk}\epsilon_{imn}$ (Paso a Paso)}
Fijemos mentalmente el primer índice en un valor concreto, por ejemplo $\mathbf{i = 1}$, y analicemos el producto $\epsilon_{1jk}\epsilon_{1mn}$ (sin sumar aún):
\begin{enumerate}
    \item Para que el primer factor $\epsilon_{1jk}$ no se anule, los índices $(j, k)$ no pueden repetir el $1$ ni repetirse entre sí. Por lo tanto, la pareja $(j, k)$ \textbf{solo puede ser $(2, 3)$ o $(3, 2)$}.
    \item Exactamente por la misma razón, para que $\epsilon_{1mn}$ no se anule, la pareja $(m, n)$ \textbf{también tiene que ser $(2, 3)$ o $(3, 2)$}.
    \item En consecuencia, el producto $\epsilon_{1jk}\epsilon_{1mn}$ es cero casi siempre, salvo en \textbf{dos únicos casos}:
    \begin{itemize}
        \item \textbf{Caso A (Emparejamiento directo: $j = m$ y $k = n$):}\\
        Ambas parejas son idénticas: o ambas son $(2, 3)$ o ambas son $(3, 2)$. En ambos escenarios los dos $\epsilon$ tienen el mismo signo, así que su producto vale $(+1)(+1) = (-1)(-1) = \mathbf{+1}$.
        \item \textbf{Caso B (Emparejamiento cruzado: $j = n$ y $k = m$):}\\
        Una pareja es $(2, 3)$ y la otra está invertida como $(3, 2)$. Un factor vale $+1$ y el otro vale $-1$, por lo que su producto vale $(+1)(-1) = \mathbf{-1}$.
    \end{itemize}
\end{enumerate}
\textbf{Pregunta clave:} ¿Qué combinación de deltas de Kronecker vale $+1$ cuando $(j=m, k=n)$ y vale $-1$ cuando $(j=n, k=m)$?
\[
\mathbf{\delta_{jm}\delta_{kn} - \delta_{jn}\delta_{km}}.
\]
*(Comprobación adicional: si $j=k$, la expresión da $\delta_{jm}\delta_{jn} - \delta_{jn}\delta_{jm} = 0$, coincidiendo con $\epsilon_{i jj} = 0$. Al sumar sobre los tres valores posibles de $i \in \{1, 2, 3\}$, para cada pareja de índices distintos $\{j, k\}$ existe un único valor de $i$ que completa la terna $\{1, 2, 3\}$, por lo que la fórmula se mantiene exacta).*
\end{intuitionbox}

\begin{theorembox}{La Identidad Maestra $\epsilon-\delta$}
En $\R^3$, la contracción de dos símbolos de Levi-Civita sobre un índice común satisface:
\[
\mathbf{\epsilon_{ijk}\epsilon_{imn} = \delta_{jm}\delta_{kn} - \delta_{jn}\delta_{km}}.
\]
\textit{Lectura visual:} «Primero con primero y segundo con segundo ($\delta_{jm}\delta_{kn}$), MENOS primero con segundo y segundo con primero ($\delta_{jn}\delta_{km}$)».
\end{theorembox}

% =====================================================================
\manualsection[5. La Recompensa]{5. La Recompensa: Demolición en 3 Líneas y Aplicaciones Físicas}

Ahora que hemos construido nuestras herramientas desde los cimientos geométricos, resolvamos en tres renglones el problema que nos tomó dos páginas en la Sección 1.

\begin{examplebox}{Demolición de la Identidad BAC--CAB en 3 Renglones}
\begin{align*}
[\vec{A} \times (\vec{B} \times \vec{C})]_i &= \epsilon_{ijk} A_j (\epsilon_{kmn} B_m C_n) = \epsilon_{kij}\epsilon_{kmn} A_j B_m C_n && \text{(1. Producto cruz en índices)} \\
&= (\delta_{im}\delta_{jn} - \delta_{in}\delta_{jm}) A_j B_m C_n && \text{(2. Identidad } \epsilon-\delta \text{)} \\
&= B_i (A_j C_j) - C_i (A_j B_j) = [\vec{B}(\vec{A}\cdot\vec{C}) - \vec{C}(\vec{A}\cdot\vec{B})]_i \quad \blacksquare && \text{(3. Filtro de la delta)}
\end{align*}
¡Observa por qué en la Sección 1 tuvimos que sumar y restar $A_x B_x C_x$! En la notación de índices, el término $j=i$ se cancela automáticamente entre $\delta_{im}\delta_{jn}$ y $-\delta_{in}\delta_{jm}$ sin que tengamos que hacer ningún truco manual.
\end{examplebox}

\manualsubsection{El Principio de Simetría vs. Antisimetría en la Física}
Cuando contraemos un objeto simétrico en dos índices ($S_{ij} = S_{ji}$, como el producto de componentes de un mismo vector $v_i v_j$) con un objeto antisimétrico ($A_{ij} = -A_{ji}$, como $\epsilon_{ijk}$), los términos con $i \neq j$ vienen en parejas que se cancelan exactamente:
\[
S_{12} A_{12} + S_{21} A_{21} = S_{12} A_{12} + S_{12}(-A_{12}) = 0.
\]
Formalmente, renombrando los índices mudos $i \leftrightarrow j$:
\[
S_{ij} A_{ij} = S_{ji} A_{ji} = (S_{ij})(-A_{ij}) = -S_{ij} A_{ij} \implies \mathbf{S_{ij} A_{ij} = 0}.
\]

\begin{examplebox}{Aplicación Física Inmediata: Fuerza Magnética de Lorentz}
La fuerza magnética sobre una partícula de carga $q$ es $\vec{F} = q(\vec{v} \times \vec{B})$, o en índices $F_i = q\,\epsilon_{ijk} v_j B_k$. La potencia mecánica (trabajo por unidad de tiempo) es:
\[
P = \vec{F}\cdot\vec{v} = F_i v_i = q\,\epsilon_{ijk} (v_i v_j) B_k.
\]
Como $v_i v_j$ es simétrico bajo $i \leftrightarrow j$ y $\epsilon_{ijk}$ es antisimétrico bajo $i \leftrightarrow j$, su contracción es **idénticamente cero**:
\[
\mathbf{P = \frac{dE_{\text{cin}}}{dt} = 0}.
\]
Un campo magnético estático jamás puede acelerar la rapidez de una partícula ni realizar trabajo mecánico.
\end{examplebox}

\manualsubsection{El Determinante como Volumen Orientado en Índices}
En la Semana 1 vimos que el triple producto escalar $\vec{u}\cdot(\vec{v}\times\vec{w})$ mide el volumen orientado del paralelepípedo formado por tres vectores. En notación de índices, este volumen es:
\[
\text{Vol}(\vec{u}, \vec{v}, \vec{w}) = \vec{u}\cdot(\vec{v}\times\vec{w}) = u_i (\vec{v}\times\vec{w})_i = \mathbf{\epsilon_{ijk} u_i v_j w_k}.
\]
Si ahora agrupamos los tres vectores como las tres columnas de una matriz $3 \times 3$, digamos $M_{i1} = u_i$, $M_{j2} = v_j$ y $M_{k3} = w_k$, el determinante de la matriz $M$ es precisamente el volumen orientado de sus columnas:
\[
\mathbf{\det(M) = \epsilon_{ijk} M_{i1} M_{j2} M_{k3}}.
\]
¿Y qué ocurre si permutamos el orden de las tres columnas tomando índices arbitrarios $(a, b, c)$ en lugar de $(1, 2, 3)$? Intercambiar dos columnas invierte la orientación de la caja y cambia el signo del volumen exactamente igual que $\epsilon_{abc}$. Por lo tanto, descubrimos la identidad fundamental del determinante:
\begin{theorembox}{Identidad Tensorial del Determinante}
Para cualquier matriz $M$ de $3 \times 3$:
\[
\mathbf{\epsilon_{ijk} M_{ia} M_{jb} M_{kc} = \det(M)\,\epsilon_{abc}}.
\]
\end{theorembox}
*(Guarda esta joya en la memoria: en la Semana 3 nos permitirá demostrar en un solo renglón por qué el producto cruz es invariante bajo rotaciones propias con $\det(R) = +1$, y en la Semana 4 explicará el factor Jacobiano en coordenadas curvilíneas).*

\end{document}
